Der Schweredruck nimmt linear mit der Tiefe zu, $\Delta p = \rho\,g\,\Delta h$. Aufgelöst nach der Dichte:
\begin{align}
 \rho &= \rhoFF \\
   &= \frac{\dpa}{\gE\cdot\hab} \\
   &= \rhoF \approx \result{\rhoFP{0}{2}}
\end{align}
Nach der Dichtetabelle (FoTaBe S.~187) hat nur \result{\text{Quecksilber}} eine so grosse Dichte ($\rho_\text{Hg} = \rhoHgO$).
